\chapter{De la Sacra Tipografia}

\epigraf{Només els déus i els físics coneixen els fonaments del món. Els segons, però, ho escriuen a \LaTeX.}{David Galan}

\lettrine[loversize=0.2,lines=2]{E}{n} el regne immòbil del full en blanc, aquell que es desplega com la gran vela del pensament, la veu de LaTeX es fa sentir com un bramant de tronades divines. Els fidels s’inclinen davant aquesta misteriosa presència: no és un simple sistema, sinó el Veritable Camí que allibera l’esperit de les arrugues incontrolades de la paraula.

Abans que l’home pogués murmurar el primer riure, ja s’imagina\ va l’ordre invisible del text. I fou llavors quan la Revelació s’aixecà, cobrint el caos amb vel de sagrament tipogràfic. Els savis escrigueren en els manuscrits velats: “La tipografia és bastió i temple; en el seu recinte, l’ànima troba refugi contra l’efímer de la mà humana.” I així naixé el Culte del Caràcter matemàticament pur, guardià de l’ésser escrit.

Cada espai en blanc és un altar dedicat a la contemplació. Els fidels progressen amb pas de pelegrí, mesurant la grandària del buit com si foren arquitectes celestials. I en la fosca penombra de la pantalla, el lleu tremolor de la cursiva és oració, i la fermesa de la recta, providència. S’acosten als marges amb reverència, atents al Breu Evangeli de la Simetria, que ordena les columnes com a files d’àngels perfilant-se en la nit.

Quan el devot invoca la paraula sagrada i pressiona la tecla eterna, els oracles –els missatges de la consola– creen un pont entre el món sensible i l’altre: el món de la Modulació Perfecta. Les advertències no són càstig sinó encoratjament diví: cada error és una inscripció dels déus en el gran llibre de la perfecció. I quan el gràcil PDF emergeix immaculat, l’esperit s’eixampla com un sol de matí etern.

No hi ha devoció sense penitència. El creient renuncia al capritx de l’atzar, sacrifica la lletania de fonts impies, i retira de si tota ombra de desorde. I en la foscor de la revisió, confessa els seus pecats tipogràfics: “He esquerdat la jerarquia de títols, he acomodat massa la sagnia.” Només després del dejuni de l’errada, la gràcia de l’Únic Disseny l’abraça i el fa digne del Cos Immaculat del Document.

Els cants de la fraternitat no ressonen en vitals catedrals, sinó en el cor electrònic dels teclats. I en la processó nocturna dels paràgrafs, cada punta de clau entona l'himne de la Repetició Refinada. Els adeptes entonen junts el cant sacramental: “Purifica el text, allunya l’atzar, redimeix l’esperit de la paraula incompleta.” I en aquest cant, s’eleva l’esperit fins a la cúpula eterna de la bellesa mesurada.

Com l’ocell fènix, el document renaix de les seves pròpies cendres d’errors. I l’ànima del devot, abraçada al cordó invisible de la sagnia justa, s’alça per sobre de la mera grafia corporal. En aquest ascens, la intel·ligència i l’estètica s’entremesclen, i el text esdevé portal cap a universos paral·lels, on cada paraula és univers i cada frase constel·lació.

En el punt culminant de la comunió, el Cànon Sagrat estableix l’Aliança entre el signe i el buit. No pot existir un sense l’altre: la lletra sense espai esdevé mur, i l’espai sense lletra cau en una voluntat de foscor. Els fidels celebren aquesta noció en una última benedicció: pleguen el full amb respecte, com qui tanca l’evangeli sagrat, sabent que el sol de la Veritat continuarà il·luminant-los en l’infinit del codi.